%%%%%%%%%%%%%%%%%%%%%%%%%%%%%%%%%%%%%%%%%%%%%%%%%%%
%
%  CHAPTER IV -- COMPLETED RESEARCH II:
%  VOLUMETRIC BRAVAIS WEAVES
%
%%%%%%%%%%%%%%%%%%%%%%%%%%%%%%%%%%%%%%%%%%%%%%%%%%%

\chapter{\uppercase{Completed Research II: Volumetric Bravais Weaves --- A Lattice-Based Framework for Discovery, Design, and Classification}}
\label{cha:volumetric-bravais}

This chapter summarizes the core theoretical contribution of the dissertation, developed across a constructive framework paper~\cite{yildiz2025volumetric} and a companion enumeration paper~\cite{akpanya2025enumeration}, which together define, construct, and count \emph{volumetric Bravais weaves}.

\section{From Tessellations to Weaves: Bravais Lattices and Wigner--Seitz Cells}
\label{sec:bravais}

Chapter~\ref{cha:preliminaries} introduced cell-induced topological elements informally, relative to an arbitrary space-filling 3D tessellation. To make this construction systematic, we restrict attention to tessellations generated by \emph{Bravais lattices}: infinite sets of points related to one another by integer combinations of a fixed set of basis vectors, characterized by the lengths of those basis vectors and the angles between them. Each Bravais lattice point induces a \emph{Wigner--Seitz cell} -- the region of space closer to that lattice point than to any other -- and the collection of these cells tiles space with a single, repeated cell shape. This is precisely the key insight of the framework: a tessellation can be turned into a weave by immersing one-dimensional thread elements into the cells, and because every Wigner--Seitz cell of a given Bravais lattice is congruent and has the same local symmetry, the cell's own point-group symmetry determines both how many distinct local thread configurations that cell type admits and how those configurations map onto one another.

Restricting to the isometries of Bravais lattices and their Wigner--Seitz cells yields exactly five three-dimensional cell types (up to combinatorial type): the cube (cubic primitive, cP), the hexagonal prism (hexagonal primitive, hP), the rhombic dodecahedron (cubic face-centered, cF), the elongated dodecahedron (orthorhombic face-centered, oF), and the truncated octahedron (cubic body-centered, cI). Each cell type admits a different -- and, in general, different-sized -- set of local thread configurations; enumerating and relating these configurations across all five cell types is the design space targeted by this dissertation, with cP treated as the fully worked example and hP, cF, oF, and cI enumerated to varying degrees of completeness (Chapter~\ref{cha:proposed-work} discusses the remaining work).

\section{Local Configurations and cP-Weaves}
\label{sec:cp-weaves}

For the cube (cP), a \emph{line configuration} is a straight thread element that enters and exits the cell through a pair of opposite (parallel) faces. Up to the symmetries of the cube, there are exactly eight such local line configurations. A \emph{cP-weave} on an $n \times n \times n$ cubic domain is then a map that assigns one of these eight configurations to each of the $n^3$ cells, subject to a compatibility rule at cell interfaces (Section~\ref{sec:symmetry}). Because each cell can independently take one of eight labels, the number of \emph{raw} assemblies for an $n\times n\times n$ domain is $8^{n^3}$ before any symmetry or compatibility constraints are applied -- already establishing that the design space grows doubly exponentially in the linear domain size $n$.

\subsection{Symmetry Analysis}
\label{sec:symmetry}

Two families of symmetry act on this raw design space:

\begin{itemize}
\item \textbf{Rotational/mirror symmetry of the cell.} The cube's automorphism group is the octahedral group, with 24 rotational symmetries (48 total, including mirror operations). Two assemblies related by one of these symmetries represent the same physical weave.
\item \textbf{Translational symmetry of the lattice.} Because the cubic array is periodic, the assembly is also invariant, up to relabeling, under lattice translations -- an additional symmetry not present for a single, isolated cell.
\end{itemize}

Combining these gives an explicit lower bound on the number of distinct $n\times n\times n$ volumetric weaves: for a domain with $48$ rotational/mirror symmetries and $n^3$ translational symmetries acting on $8^{n^3}$ raw cell-labelings, the number of \emph{orbits} (distinct weaves) is at least on the order of $8^{n^3}/(48\,n^3)$ -- still exponential in the domain size, and computable in principle by brute-force orbit enumeration for small $n$.

\subsection{Case Study: Plain-Weave cP-Structures}

Restricting to the classical \emph{plain-weave} rule -- every interface between adjacent cells must change local configuration -- substantially shrinks the admissible design space: each cell's configuration can be represented by three binary vectors of length $n$ (one per axis), giving a lower bound that is still exponential in $n$ but far smaller than the unrestricted case. Implementing an orbit-finding algorithm for plain cubic volumetric weaves under the rotational, mirror, and translational symmetries above, we find that a $2\times 2\times 2$ domain admits exactly $5$ inequivalent plain cP-weaves, and a $4 \times 4 \times 4$ domain admits exactly $20$.

\subsection{Beyond cP: hP and Other Lattices}

The same local-configuration analysis applied to the hexagonal-prism (hP) cell -- which carries four threads per cell rather than one -- yields $76$ distinct local configurations, an order-of-magnitude increase over the eight configurations of cP. This growth illustrates the paper's central claim: even restricting to a single, simple example (cP-weaves), the design space of volumetric Bravais weaves is already at least one order of magnitude larger than that of conventional two-way, two-fold planar weaves, and grows substantially further once other Wigner--Seitz cell types are included. Completing the enumeration of local configurations for cF, oF, and cI is part of the proposed work described in Chapter~\ref{cha:proposed-work}.

\section{Rigorous Enumeration: Partial Cubic Primitive Weaves}
\label{sec:enumeration}

The symmetry counting in Section~\ref{sec:symmetry} gives a lower bound; the companion paper~\cite{akpanya2025enumeration} makes this exact using classical enumeration combinatorics. There, a \emph{partial} $n\times n\times n$ cubic primitive weave is modeled as a map from the $n^3$ cells of the domain to the set of eight local configurations up to homotopy, and the group
\[
G_n \;=\; \mathrm{Rot}(C) \ltimes \mathbb{Z}_n^3
\]
(the semidirect product of the cube's rotation group with the group of translations modulo $n$ in each coordinate direction) acts naturally on the set of such maps by simultaneously rotating the ambient cube and permuting cells according to the induced translation. By computing the cycle structure and fixed-point count of every element of $G_n$ acting on (equivalence classes of) partial weaves, and applying Burnside's lemma in the de Bruijn enumeration framework, we obtain a closed-form count of inequivalent partial $n\times n\times n$ cubic primitive weaves:

\begin{theorem}[Enumeration of partial cP-weaves, informal statement; see~\cite{akpanya2025enumeration} for the full derivation]
\label{thm:enumeration-informal}
For every $n \geq 3$, the number $\psi(n)$ of inequivalent partial $n \times n \times n$ cubic primitive weaves is given by
\[
\psi(n) \;=\; \frac{1}{|G_n|} \sum_{g \in G_n} \operatorname{Fix}_{\Omega^{\Delta_n}}(g),
\]
where $\Delta_n$ is the set of $n^3$ cells, $\Omega$ is the set of equivalence classes of local configurations, and the fixed-point counts $\operatorname{Fix}_{\Omega^{\Delta_n}}(g)$ are determined explicitly by the cycle type of $g$ acting on $\Delta_n$ and on $\Omega$. The case $n=2$ is exceptional, since the rotation and translation subgroups do not form a semidirect product of order $24n^3$, and is treated separately.
\end{theorem}

This result upgrades the exponential lower bound of Section~\ref{sec:symmetry} to an exact count, and provides the same closed-form machinery (cycle-index computation for a lattice-symmetry group acting on a finite set of local labels) that Chapter~\ref{cha:proposed-work} proposes to extend to the remaining Bravais cell types (hP, cF, oF, cI) as time permits.

\section{Summary and Impact}

Taken together, these results establish volumetric Bravais weaves as a formal, dimension-consistent generalization of Gr\"unbaum and Shephard's planar isonemal fabrics: a volumetric weave is a labeling of the cells of a Wigner--Seitz tessellation by local, cell-induced topological elements, subject to the symmetry group of the underlying Bravais lattice. This framework (i) reproduces known planar weaves as the degenerate 2D case, (ii) provides an explicit, enumerable design space for cP-weaves that is provably at least an order of magnitude richer than the planar case, and (iii) is exact rather than approximate, thanks to the group-theoretic enumeration in Section~\ref{sec:enumeration}. Chapter~\ref{cha:nonwoven} extends the same cell-induced perspective to structures that relax the strict over--under interlacing rule (nonwoven textiles), and Chapter~\ref{cha:algebraic-calculus} develops algebraic machinery to \emph{design} within this space directly, rather than enumerating it after the fact.
